\section{Concetti di base}
Ribadiamo che lo scopo della programmazione multicore è mantenere la CPU quanto più attiva possibile, minimizzando i tempi in cui non sta elaborando
processi. Questo concetto, nei sistemi multicore, è esteso ad ogni core di elaborazione del sistema, la cui gestione avviene tramite
\textbf{scheduling}.\par
I processi in esecuzione si compongono di un \textbf{ciclo di elaborazione CPU} e un \textbf{ciclo di attesa di completamento operazioni I/O}.
L'esecuzione ha inizio con una sequenza di operazioni CPU, detta \textbf{CPU burst}, seguite da operazioni I/O, chiamate \textbf{I/O burst}; questi due
elementi si avvicendano fino al termine del processo con la chiamata opportuna al sistema operativo.\par
Dove le durate dei CPU burst variano in base al tipo di processo e all'architettura, in genere, i processi I/O-bound presentano CPU burst di breve
durata, mentre i CPU-bound ne mostrano di lunga durata.\par
Ogni volta che la CPU passa in stato di inattività, è compito dello \textbf{scheduler a breve termine} selezionare il processo dalla ready queue e dargli
il controllo dell'unità. Le queues non sono necessariamente code FIFO; possono essere implementate anche con logica LIFO o altre strutture dati come
alberi o liste concatenate; tuttavia, gli elementi al suo interno devono sempre essere i PCB dei processi. Le decisioni per lo scheduling sono
effettuate in quattro contesti: \begin{itemize}
    \item Quando un processo passa da esecuzione ad attesa, come per le richieste I/O oppure quando si invoca wait().
    \item Quando un processo passa da esecuzione a pronto, come quando si segnala un interrupt.
    \item Quando un processo passa da attesa a pronto, come quando si è completata un'operazione I/O.
    \item Quando un processo termina l'esecuzione.
\end{itemize}

\noindent Quando lo scheduler interviene nei casi $1,4$ si dice che è \textbf{senza prelazione} e implica che quando si assegna il controllo della CPU
a un processo, questo lo mantiene fino al momento del rilascio, dovuto al suo termine o al passaggio in stato di attesa.\par
Negli altri due casi, lo schema di scheduling è \textbf{con prelazione}; sebbene sia la strategia attuata da tutti i moderni sistemi operativi, presenta
alcuni svantaggi, come portare a race condition\footnote{Un processo potrebbe modificare dati che un altro in attesa deve leggere. Il risultato è 
non deterministico.} quando i dati sono condivisi tra i processi, dinamica valida anche per il kernel. Si potrebbe pensare quindi di scegliere uno
schema senza prelazione, ma non è adeguato per le applicazioni in tempo reale. È quindi necessario applicare meccanismi per prevenire la formazione di
race condition.\par
Un altro elemento coinvolto nello scheduling è il \textbf{dispatcher}, il modulo che passa effettivamente il controllo della CPU al processo
scelto dallo scheduler. Le mansioni comprendono il context switch da un processo all'altro, il passaggio alla user-mode e il salto alla posizione
adeguata dei programmi utente per la loro esecuzione.\par
Idealmente, i dispatcher dovrebbero essere quanto più ottimizzati possibile poiché i cambi di contesto avvengono frequentemente. Chiamiamo
\textbf{latenza di dispatch} il tempo richiesto per fermare un processo e avviarne un altro cedendo il controllo della CPU. In particolare, esistono
due tipi di context switch: \textbf{volontario}, se un processo cede il controllo della CPU in seguito alla richiesta di una risorsa attualmente non
disponibile, e \textbf{involontario} quando il comando è sottratto dal processo, per esempio quando è esercitata la prelazione.\newline

\noindent Esistono più algoritmi di scheduling con diverse caratteristiche, ed è necessario prenderle in considerazione prima di implementare quello più
adatto al contesto in esame. Consideriamo alcuni criteri per una scelta consapevole: \begin{itemize}
    \item \textbf{Utilizzo della CPU}: L'unità deve essere quanto più attiva possibile e il suo utilizzo si misura in percentuale. In un sistema a basso
    carico è di circa $40\%$, mentre se è ad alto carico, si aggira sul $90\%$.
    \item \textbf{Throughput}: Numero di processi completati nell'unità di tempo.
    \item \textbf{Tempo di completamento}: Intervallo di tempo che intercorre tra sottomissione del processo al suo completamento. È il risultato della
    somma dei tempi di attesa per l'inserimento in memoria, nella ready queue, durante l'esecuzione e nelle operazioni di I/O.
    \item \textbf{Tempo di attesa}: Somma dei tempi passati nella ready queue.
    \item \textbf{Tempo di risposta}: Tempo necessario per iniziare la risposta ad una richiesta.
\end{itemize}
\noindent Generalmente, è una buona prassi cercare di massimizzare i tempi di utilizzo della CPU e minimizzare quelli di attesa, risposta e completamento.
In particolare, con sistemi interattivi come quelli desktop, è auspicabile ridurre al minimo la varianza del tempo di risposta per garantire
un'esperienza stabile ed efficace.

%

\section{Algoritmi di scheluding}

%

\section{Scheduling dei thread}

%

\section{Scheduling per sistemi multiprocessore}

%

\section{Scheduling real-time}

%

\section{Valutazione algoritmi}